\documentclass[UTF8,a4paper,11pt,fontset=fandol]{ctexart}
\usepackage[margin=25mm]{geometry}
\usepackage{amsmath,amssymb,amsthm,mathtools,bm}
\usepackage{booktabs,longtable,array,enumitem}
\usepackage[numbers,sort&compress]{natbib}
\usepackage{xurl}
\usepackage[colorlinks=true,linkcolor=blue,citecolor=blue,urlcolor=blue]{hyperref}
\usepackage{microtype}
\providecommand{\doi}[1]{\href{https://doi.org/#1}{\nolinkurl{#1}}}
\setlength{\emergencystretch}{3em}
\setlength{\parskip}{0.25em}
\setlist{nosep}
\numberwithin{equation}{section}
\newtheorem{theorem}{定理}[section]
\newtheorem{proposition}[theorem]{命题}
\newtheorem{lemma}[theorem]{引理}
\newtheorem{corollary}[theorem]{推论}
\theoremstyle{definition}
\newtheorem{definition}[theorem]{定义}
\newtheorem{assumption}[theorem]{假设}
\theoremstyle{remark}
\newtheorem{remark}[theorem]{说明}
\newtheorem{counterexample}[theorem]{反例}
\DeclareMathOperator{\Fix}{Fix}
\DeclareMathOperator{\Stab}{Stab}
\DeclareMathOperator{\rank}{rank}
\DeclareMathOperator{\tr}{tr}
\DeclareMathOperator{\im}{im}
\DeclareMathOperator{\diag}{diag}
\DeclareMathOperator{\Span}{span}
\newcommand{\R}{\mathbb R}
\newcommand{\C}{\mathbb C}
\newcommand{\Z}{\mathbb Z}
\newcommand{\Q}{\mathbb Q}
\newcommand{\T}{\mathsf T}
\newcommand{\I}{\mathrm I}
\newcommand{\eps}{\varepsilon}
\newcommand{\norm}[1]{\left\lVert #1\right\rVert}
\newcommand{\code}[1]{\texttt{\detokenize{#1}}}
\newcommand{\scope}{本文是项目技术论证，未经独立同行评审。定理证明针对明列的数学假设；
它们不自动证明某个二进制、数据库、浮点容差或官网版本满足这些假设。}

\title{ISODISTORT 计算能力的理论基础：\\空间群、序参量、模式分解与可验证性}
\author{CRIS 项目理论资料整理}
\date{2026年10月15日}
\begin{document}
\maketitle
\begin{abstract}
本文以官方帮助列明的功能和原始方法论文为边界，建立 ISODISTORT 计算能力的统一数学框架。
讨论结构对称性、Wyckoff 轨道、k-star 与物理不可约表示、各向同性子群、
四种搜索/分解方法、位移/占位/磁/转动/应变、超空间、等价畴、子群树、
结构生成、衍射和模式文件交换。对可在本文内封闭证明的核心合同给出推导；
对依赖已发表分类算法、表格或具体软件约定的环节列出文献前提与审查条件。
特别区分对称允许、能量稳定和实现正确，避免把一份理论说明误用为整个网站的形式化认证。
\end{abstract}
\noindent\textbf{关键词：}ISODISTORT；ISOTROPY；晶体学；表示论；各向同性子群；超空间；理论验证。
\par\scope
\clearpage
{\small\tableofcontents}
\clearpage

\section{范围、方法与文献来源}\label{sec:scope}
ISODISTORT 是 ISOTROPY 算法的交互界面；其起始方法论文是
Campbell、Stokes、Tanner 与 Hatch 的 ISODISPLACE 论文
\citep{Campbell2006}。这里“全部计算能力”指本次保存的官方帮助快照
\citep{OfficialHelp} 中列出的数学能力，而不指未知未来版本或所有 ISOTROPY 兄弟程序。
本地 CRIS 的产品范围更小：当前网页提供 Method 1--4 及相应导出，
官网的内置可视化、Generate、Domains、完整子群树等不能因本文收录理论而视为本地已实现。

本文采用原始论文、作者官方文档和 IUCr 标准；独立交叉来源包括
AMPLIMODES 与模式晶体学综述\citep{Amplimodes2009,PerezMato2010}、
磁超空间综述\citep{Magnetic2012}和 spglib 方法论文\citep{Spglib2024}。
International Tables A、A1、D 的出版入口用于规定权威参考范围
\citep{ITA,ITA1,ITD}；本次未获得全部章节正文，不伪造其页码或声称逐页核验。
检索平台只用于发现，最终引用以论文或官方页面为准。
所有下载具有 URL、获取时间、字节数与 SHA-256；扫描件和仅有摘要的来源分别标记。

\paragraph{证明对象。}
本文证明数学映射、判定条件和不变量；文献中的算法完备性只在其原始假设下使用。
不同软件可能使用不同表示基、原点、行列约定和归一化。
因此“来自同一理论”不能推出“逐列输出相同”，更不能推出每个数值分支无误。
本地实现特有的补充证明和未认证项另见
\href{unsupported_algorithm_proofs.pdf}{《数学与晶体学补充论证》}。

\section{坐标、晶格与仿射群}\label{sec:affine}
取直接基 $A=(a_1,a_2,a_3)$，$r=Ax$；度量 $G=A^\T A>0$。
倒易基使用 $A^{-\T}$，在 Fourier 相位中显式保留 $2\pi$。
行基约定 $P=A^\T$ 必须在接口处转置，不能只改文字说明。

\begin{proposition}[仿射复合与换基]\label{prop:affine}
$\{R|t\}\{S|u\}=\{RS|Ru+t\}$。
若 $x=Bx'+p$，则同一操作在新坐标中为
\[
 R'=B^{-1}RB,\qquad t'=B^{-1}(Rp+t-p).
\]
度量变为 $G'=B^\T GB$，倒易分数坐标变为 $k'=B^\T k$。
\end{proposition}
\begin{proof}
逐次代入 $x\mapsto Sx+u\mapsto R(Sx+u)+t$ 得复合律。
再代入 $x=Bx'+p$ 并左乘 $B^{-1}$ 得 $R',t'$。
由 $\norm{ABx'}^2=x'^\T B^\T GBx'$ 得度量，
由 $k^\T Bx'=(B^\T k)^\T x'$ 得倒易坐标。
\end{proof}
这是所有 basis/origin 处理的基本合同\citep{Hatch1985,ITA}。
正交旋转条件是 $R^\T GR=G$，一般不是 $R^\T R=\I$。
标准 setting、Hall 标识和原点选择的数据库对照不能用空间群号相同代替。

\begin{proposition}[超胞指数与中心化]
对原胞基上的整数非奇异 $S$，
$[\,\Z^3:S\Z^3\,]=|\det S|$。
若使用中心化惯用胞，则必须先换到原胞再求指数。
\end{proposition}
\begin{proof}
整数伴随矩阵恒等式保证 $(\det S)\Z^3\subseteq S\Z^3$，故商群有限，记阶为 $N$。
$N$ 个整数余类代表各加上 $[0,1)^3$，给出体积为 $N$ 的子格基本域；
$S[0,1)^3$ 则给出体积为 $|\det S|$ 的基本域。
同一格的两个有界基本域由平移交集分割互相重组，体积相同，故 $N=|\det S|$。
惯用胞不生成全部平移格，故只用惯用胞整数三元组会遗漏中心化平移。
\end{proof}
计算原子重数、倒易格、畴数和 $A_s/A_p$ 时，这个区别具有实际影响。
原胞、中心化惯用胞和倒易格的定义亦见 International Tables C 的公式章节
\citep{ITCGeometry}。

\section{结构对称性、Wyckoff 轨道与输入识别}\label{sec:symmetry}
设每个原胞内的原子为带物种、占位等标签的有限集合。
精确对称操作必须保持度量，并使各物理标签按指定张量类型变换。
对占位/物种不变的结构，条件为存在双射 $\pi$ 及整数 $n_j$，
\[
 Rx_j+t=x_{\pi(j)}+n_j.
\]
\begin{proposition}[理想输入下的有限对称搜索]
在固定原胞基中，对 $R\in GL(3,\Z)$ 要求 $R^\T GR=G$，
可能的 $R$ 有限。每个 $R$ 的平移候选模 $\Z^3$ 可从一个选定原子到相容原子的映射产生。
完整检查所有候选双射即可得到精确空间群操作模平移格的集合。
\end{proposition}
\begin{proof}
$R$ 的第 $j$ 列 $r_j$ 满足
$\lambda_{\min}(G)\norm{r_j}^2\leq r_j^\T Gr_j=G_{jj}$，
所以各列在有限整数盒内。
对固定原子 $x_0$，任一有效 $t$ 必须等于 $x_i-Rx_0$ 模整数；
相容 $i$ 有限。检验全部原子是充要条件。
\end{proof}
FINDSYM 与 spglib 提供实际识别算法\citep{Findsym2005,Spglib2024}。
上述证明针对精确数据，不是对带噪结构唯一空间群的证明。
容差可能改变对称群；解析失败、部分占位以及原子匹配并列必须明确处理。

\begin{theorem}[轨道与局部对称]
有限商群作用下，位置 $x$ 的轨道大小等于 $|G|/|G_x|$，
其中 $G_x$ 是位置稳定子。限制到子群 $H$ 后，母轨道拆成互不相交的 $H$ 轨道。
\end{theorem}
\begin{proof}
映射 $gG_x\mapsto gx$ 良定义且双射，给出轨道大小。
$H$ 的“可达”关系为等价关系，故将母轨道分割成子轨道。
\end{proof}
这解释 Wyckoff 分裂与独立原子参数数目的变化。
不同标签体系中相同字母未必表示同一位置，必须比较实际 setting 和轨道。

\section{波矢、k-star、小群与实物理表示}\label{sec:k}
对母平移 $n$，Bloch 模式的特征标为
$\chi_k(n)=e^{2\pi i k^\T n}$。
两个 $k$ 等价当且仅当其差属于实际倒易格。
\begin{proposition}[倒易作用与小群]
若主动旋转的分数矩阵为 $R$，则波矢变为 $R^{-\T}k$。
小群由满足 $R^{-\T}k\equiv k$ 的操作组成，k-star 是该作用的轨道。
\end{proposition}
\begin{proof}
旋转后的标量波在 $x$ 取值为
$e^{2\pi i k^\T R^{-1}x}=e^{2\pi i(R^{-\T}k)^\T x}$。
小群是该余类的稳定子，轨道--稳定子关系给出 star 的 arm 数。
\end{proof}
非对称形的小群还含分数平移，其相位不可从点群矩阵中删除。
任意 irrep 表必须与该基、原点及 Fourier 符号匹配；
ISO-IR 与 ISO-KOV 是表示及标签关系的官方入口
\citep{Stokes2013,IsoIr,IsoKov}。

\begin{proposition}[复表示的实化]
若物理场实值，则 Fourier 系数满足 $u_{-k}=\overline{u_k}$。
对复矩阵 $D=X+iY$，其在实坐标 $(\Re u,\Im u)$ 上的作用为
$\left(\begin{smallmatrix}X&-Y\\Y&X\end{smallmatrix}\right)$，
并保持表示乘法。
\end{proposition}
\begin{proof}
展开 $Du=(X+iY)(a+ib)$，实部为 $Xa-Yb$，虚部为 $Ya+Xb$；
复矩阵乘法与该块矩阵乘法逐项相同。
\end{proof}
该实化不一定在实数域不可约；应进一步区分实型、复型和四元数型，
不能把所有复表一律当作独立的双倍物理自由度。
$k\equiv-k$ 的自共轭情形还可能使 sine 模式消失。
磁静态序参量在这里是经典实张量，不把它与量子态的反幺正自旋表示混同。

\section{五类物理变量的表示与模式投影}\label{sec:physical}
每个母原子轨道上的物理变量都形成表示空间。
若 $g$ 将站点 $j$ 送至 $\pi_g(j)$，笛卡尔旋转为 $Q_g=AR_gA^{-1}$，
则占位标量按置换变换，位移按 $Q_g$ 变换。
轴矢量按 $\det(Q_g)Q_g$ 变换。
磁矩还乘时间反演符号 $\theta_g=\pm1$，转动轴矢量则为时间偶。

\begin{proposition}[极矢量与轴矢量区别]
对正交 $Q$，$(Qa)\times(Qb)=\det(Q)Q(a\times b)$。
因此磁矩和转动与位移在反演下具有不同变换律。
\end{proposition}
\begin{proof}
对任意 $c$，左式与 $Qc$ 的点积等于有向体积
$\det(Q)\det[a,b,c]$；右式与 $Qc$ 的点积相同。
由于 $Q$ 可逆、$c$ 任意，两个向量相等。
\end{proof}
含时间反演的灰群和磁群需采用相应表示与标准 setting
\citep{Magnetic2012,IsoMag}；本地菜单出现“magnetic”不意味着全部磁超空间输入已验收。

\begin{theorem}[对称内积与完全可约性]
有限群实/复表示可装备不变正定内积；
任一不变子空间的正交补仍不变。
\end{theorem}
\begin{proof}
从任一 $W_0>0$ 构造
$W=\sum_{g\in G}D(g)^*W_0D(g)>0$，右乘重排给 $D(h)^*WD(h)=W$。
若 $v\perp U$ 且 $U$ 不变，则
$\langle D(h)v,u\rangle=\langle v,D(h)^{-1}u\rangle=0$。
逐次取不可约子空间及正交补便得直和分解。
\end{proof}

\begin{theorem}[不可约分量投影]\label{thm:irproj}
设 $D$ 为有限群的酉表示，$\alpha$ 为维数 $d_\alpha$ 的复不可约表示。
\[
 P_\alpha=\frac{d_\alpha}{|G|}
 \sum_{g\in G}\overline{\chi_\alpha(g)}\,D(g)
\]
是对应等型分量的投影，且重数为
$m_\alpha=|G|^{-1}\sum_g\overline{\chi_\alpha(g)}\chi_D(g)$。
\end{theorem}
\begin{proof}
先证明 Schur 正交性：对任意矩阵 $X$，平均
$|G|^{-1}\sum_g D_\beta(g)X D_\alpha(g)^{-1}$ 是 intertwiner。
不同不可约表示间它为零；相同表示间由 Schur 引理等于
$\tr(X)\I/d_\alpha$。Schur 引理本身由非零 intertwiner 的核、像不变，
及复数域中取一个特征值后考察 $T-\lambda\I$ 的核推出。
令 $X$ 依次取矩阵单位，获得矩阵元正交性及字符正交性。
在完全可约分解上逐块应用，$P_\alpha$ 在 $\alpha$ 块为单位，
在其余块为零；取迹即得重数式。
\end{proof}
在实表示中应先正确复化/实化，不用复字符公式错误计算物理维数。
投影每个站点局部试探矢量，再去除线性相关项，可构造整体模式；
其原始算法见 \citet{Stokes1991}。
重复 irrep 实例中的基可以继续正交旋转，模式标签不是纯由空间群号决定的。

\paragraph{占位与应变的附加约束。}
占位概率须满足 $0\leq o_j\leq1$，固定组分还要求加权总占位守恒。
模式线性空间本身不保证这些不等式。
笛卡尔小应变是对称二阶张量，作用为 $E\mapsto QEQ^\T$，
空间维数为六。纯平移在均匀应变上恒等，故仅靠应变不能破坏平移对称。

\paragraph{刚体转动的线性/有限区别。}
无穷小位移 $\delta r=\omega\times(r-r_0)$ 描述切空间。
有限转动应为 $r'=r_0+\exp([\omega]_\times)(r-r_0)$；
反对称矩阵的指数正交，所以每个刚体内部长度保持。
但相邻刚体共享原子的相容条件仍须同时满足；
线性零模不必积分为有限的无应变机构。
转动模式的实际参数化见 \citet{Muller2014,Howard1998,HowardErratum2002}。
这种区别也防止把 ISODISTORT 的转动变量与 ISOTILT 的全部刚性网络算法混为一谈。

\section{序参量、各向同性子群及等价性}\label{sec:isotropy}
给定表示 $D$ 与序参量 $\eta$，各向同性子群定义为
$G_\eta=\{g:D(g)\eta=\eta\}$。
\begin{proposition}
$G_\eta$ 是子群；若 $\eta'=D(a)\eta$，
则 $G_{\eta'}=aG_\eta a^{-1}$。
\end{proposition}
\begin{proof}
单位元固定 $\eta$；若 $g,h$ 固定 $\eta$，
则 $gh$ 和 $g^{-1}$ 也固定。再由
$D(g)D(a)\eta=D(a)\eta\iff D(a^{-1}ga)\eta=\eta$ 得共轭关系。
\end{proof}
等价性必须说明允许的共轭/基变换群；
不同嵌入可能有同一空间群类型而不是同一畸变。
原始等价分类见 \citet{Stokes1985}。

\begin{theorem}[有限表示下的完整但低效的枚举]
给定有限 $G$ 的完整乘法表与表示，可枚举其所有子集，
保留子群，对每个 $H$ 求 $\Fix(H)$ 并应用精确稳定子判据，
再按所声明的等价关系去重，从而得到全部非零各向同性类型。
\end{theorem}
\begin{proof}
$G$ 有有限子集；子群条件可有限检查。任意非零 $\eta$ 的稳定子必在枚举中，
并通过固定空间与精确性检查，故无遗漏。
通过检查的每个 $H$ 由补充文档的见证定理实现，故无伪候选。
最后精确等价关系只合并同类，保留所有类别。
\end{proof}
这证明能力的数学可实现性，不声称是网站采用的高效实现。
实际算法通过平移、点群、生成元和 OPD 分层缩小搜索
\citep{Stokes2016,Stokes2017}。查表的完备性仍依赖表项本身被独立核验。

\section{四种方法的输入、输出与正确性条件}\label{sec:methods}
\subsection{Method 1：特殊 k 点上的单 irrep 搜索}
官方合同是选定特殊 k 集及单个 irrep 产生的各向同性子群，
并施加晶系、群类型、子格及最大子群等过滤\citep{OfficialHelp}。
数学路径是：有限特殊 k 表 $\to$ 物理可实现 irrep $\to$
OPD/子群 $\to$ 精确过滤。
若每一步候选集合完整、过滤等价关系正确，则集合交运算保证相对这一输入域的完备性。
特殊 k 列表不是整个连续 Brillouin 区；单 irrep 也不能代替耦合 irrep。
父点群旋转下等价子格须按实际平移格判等，不能仅比较三个显示向量。

\subsection{Method 2：指定 k 与 irrep、允许耦合和调制}
选择 k、irrep 和 OPD 后，计算其稳定子，再在各物理变量空间中求该子群的固定空间。
耦合空间必须用直和表示；子群为各分量稳定子的交。
单靠 irrep 名称相同不能认证相位原点和列身份。
公度情形可下降到有限平移商群；非公度情形需第 \ref{sec:ssg} 节的超空间理论。
任意 OPD 的数值采样只可发现某些稳定子，不证明所有特殊方向已找到。
文献算法及官方文档均区分预计算表与按需生成
\citep{Stokes2016,Stokes2017,OfficialHelp}。

\subsection{Method 3：指定点群/空间群类型与超胞的逆问题}
先固定子平移格 $L'$。若其不是母群正规子格，
不能直接把 $G/L'$ 当作群；可先取有限指数正规核
$L_0=\bigcap_{g\in G}gL'g^{-1}$，在 $G/L_0$ 中工作。
点群有限，所以该交只需有限不同项。
候选嵌入 $H$ 须保留群类型、basis、origin 和物理表示，而非只有空间群号。
通过平移相位枚举相容 k-star，再用限制表示的固定空间确定初级/次级序参量
\citep{Hatch2002,COPL2002}。
非零固定空间仅说明某些畸变保持 $H$；精确稳定子检查排除实际对称性更高者。
对由有限指数超胞约束产生的有限商群，枚举全部子群给出原则上的完备算法；
实际搜索的指数上界、超时、候选裁剪必须作为额外边界记录。
若目标嵌入与指定物理类型不相容，正确结果可能是空集。

\subsection{Method 4：给定子相的模式分解}
必须先指定母子 setting、超胞、原点和物种双射，剥离所选均匀应变规范，
得到公共参考结构上的位移 $u$。
模式矩阵 $B$ 的列必须覆盖被分解的自由参数空间。
其带度量最小二乘解见补充论证；若 $B$ 为完整独立基则残差为零，
否则残差只表示相对于所选子空间的不可表示分量。
原点/匹配不唯一、缺陷或占位变更会改变这个问题本身。
因此先做最小二乘再给晶体学身份是不充分的
\citep{Campbell2006,PerezMato2010,Amplimodes2009}。

\section{模式幅度、单位、重数与应变}\label{sec:amplitudes}
原始模式列 $b_i$、归一化列 $\hat b_i=n_i b_i$ 和物理幅度 $a_i$ 满足
$u=\sum_i a_i\hat b_i$。
定义范数时必须包含全部对称生成的原子、正确度量及约定的 primitive supercell。
由 $A_s^2$ 为总平方范数得到
$A_p=A_s\sqrt{V_p/V_s}$，体积是参考 primitive cell 的体积
\citep{OfficialHelp}。
在一组正交归一模式中，$\norm u^2=\sum_i a_i^2$；
若模式非正交，正确表达式是 $a^\T(B^\T WB)a$，不能直接取系数平方和。

\begin{proposition}[协方差传播的条件]
当 $a=Lu$ 为固定线性映射且 $\operatorname{Cov}(u)=\Sigma$ 时，
$\operatorname{Cov}(a)=L\Sigma L^\T$。
\end{proposition}
\begin{proof}
将 $a-\mathbb Ea=L(u-\mathbb Eu)$ 代入协方差定义即可。
\end{proof}
这是线性误差传播的证明，不是网站当前计算 ESD 的证明。
官方帮助的旧叙述提到 ESD，版本记录却记载 5.5 版本移除 ESD；
本库将该项标为\textbf{文档冲突/版本依赖}\citep{OfficialHelp,OfficialVersions}。

有限应变要与笛卡尔小应变、Green--Lagrange 应变和分数基对称乘子区分。
代码中的正定度量匹配方程有严格唯一解，但其与官网规范的一致性需要额外证据。
宏观应变模式不适用微观原胞重复计数；
工程剪切的因子 $2$ 与 tensor shear、Voigt 顺序须逐接口记录。

\section{非公度调制、谐波及超空间}\label{sec:ssg}
设 $d$ 个调制波矢组成 $Q\in\R^{d\times3}$，
内坐标为 $\xi=Qx+\phi\pmod{\Z^d}$。
位移/占位等周期函数可写为有限截断或收敛 Fourier 展开
\[
 f(\xi)=\sum_{m\in\Z^d}c_m e^{2\pi i m^\T\xi},
 \qquad c_{-m}=\overline{c_m}
\]
以保证实值。调制函数、卫星及超空间截面的直观解释见 \citet{Wagner2009}。
Fourier 级数在何种范数收敛取决于函数正则性；
有限谐波截断不能无条件证明对任意调制完全精确。

\begin{proposition}[截面的仿射相容性]
设超空间操作
$\widetilde g:(x,\xi)\mapsto(Rx+t,Mx+S\xi+\tau)$，
其中 $M$ 为整数 $d\times3$ 矩阵，$S\in GL(d,\Z)$。
它把固定截面 $\xi=Qx+\phi$ 映到同一截面，当且仅当
\[
 M+SQ=QR,\qquad \tau=Qt+\phi-S\phi\pmod{\Z^d}.
\]
\end{proposition}
\begin{proof}
把截面代入变换的第二分量，得到
$(M+SQ)x+S\phi+\tau$。要求它等于
$Q(Rx+t)+\phi$ 模整数。对连续 $x$ 比较线性项和常数项便得条件。
\end{proof}
一般超空间对称还可以把一个截面映到另一个相位截面；
不能将上述“固定截面”条件误用为所有非公度体相操作的限制。
坐标系和波矢规范不同会改变矩阵块的符号与位置，须按命题
\ref{prop:affine} 转换。

\begin{proposition}[内部平移的谐波表示]
内平移 $\delta$ 将第 $m$ 个复 Fourier 模式乘以
$e^{2\pi i m^\T\delta}$；
在其 cosine/sine 实子空间上这是二维正交旋转，旋转方向取决于已声明的主动/被动约定。
\end{proposition}
\begin{proof}
将 $\xi+\delta$ 代入指数函数，然后取实部与虚部，应用三角加法公式即可。
\end{proof}
对于公度截面，特征标可能自共轭或发生别名；补充论证给出阶数和秩的精确判定。
对于无理波矢，内部平移的相位群一般不是有限群，不能直接套用有限群枚举证明。

超空间群生成、标准识别、等价变换和相应 irrep 的基础由
\citet{Stokes2007,Stokes2011,Stokes2013,Smaalen2013,Stokes2017} 给出。
本次核查到 2017 论文将某些微观基函数投影问题留待后续工作，
因此不能仅凭该一篇论文声称非公度微观列定义的全部来源已覆盖。
补充论证的“非公度微观模式”一节进一步从函数拉回作用严格构造微观表示，
利用生成元的等变线性方程求出全部指定 irrep 的嵌入；这不需要把无理相位伪装成有限群。
但若站点对应、内部相位或 irrep 基约定未知，仍不能确定官网逐列身份。
本节证明仿射/Fourier 合同；完整超空间分类的证明依赖上述文献，
数据库每一条和所有谐波组合的实现仍属条件性支持。
磁超空间必须另加时间反演及磁轴矢量合同\citep{Magnetic2012}。

\section{生成结构、畴、序参量组合与子群树}\label{sec:derived}
\subsection{由幅度生成结构}
在相容模式空间中 $u=B a$。
若 $D(h)B=B$，则 $D(h)u=u$，生成结构至少保持 $H$。
特殊幅度可能恢复更高对称性；零幅度一般恢复母相。
若幅度违反占位界、造成原子重叠或进入非线性转动的错误参数域，
保持群对称仍不等于结构物理可行。

\subsection{等价畴}
\begin{theorem}[畴与陪集]
若一个畸变的精确稳定子为 $H$，其母群轨道与左陪集 $G/H$ 双射。
有限指数情形的畴数为 $[G:H]$。
\end{theorem}
\begin{proof}
$g_1\eta=g_2\eta$ 当且仅当 $g_2^{-1}g_1\in H$，
恰为相同左陪集的条件。
\end{proof}
若母平移群为 $T$，则
$[G:H]=[G:TH]\,[T:T\cap H]$。
第一个因子对应点群取向的差别，第二个对应平移/反相畴；
等式由指数乘法及 $TH/H\simeq T/(T\cap H)$ 的集合双射成立。
实际畴壁、孪晶界能量和动力学不由此给出。

\subsection{潜在主序参量与不变量}
直和序参量的稳定子取交，故可以列出达到目标 $H$ 的组合；
多个不同组合可能都足以产生同一 $H$。
对有限群和多项式 $f$，Reynolds 平均
$\mathcal R f(\eta)=|G|^{-1}\sum_g f(D(g)\eta)$ 为不变量，
且任何不变量都满足 $\mathcal R f=f$。
因此对全部固定次数单项式求平均并线性去重，可张成该次数的不变量空间。
这个证明不保证任意次数截断已经找到整个不变量环的生成元。
INVARIANTS 和耦合理论提供相关方法\citep{Invariants2003,Coupled1991}；
网站“potentially primary”不是唯一热力学驱动力的测量。
自由能系数、稳定极小值和一级/二级相变的实际判断还需要独立物理模型和实验。

\subsection{中间子群树}
当 $H\leq K\leq G$ 且 $[G:H]<\infty$，
取 $H$ 的正规核 $N=\bigcap_{g\in G}gHg^{-1}$，它具有有限指数。
中间子群对应有限群 $G/N$ 中包含 $H/N$ 的子群。
枚举它们并按包含关系取覆盖边，得到完整 Hasse 图；
若另对共轭类合并，必须记录其合并标准和丢失的嵌入重数。
这一证明解释子群树的数学对象，不认证网页资源限制下每次生成都完整。

\section{衍射与可视化的理论边界}\label{sec:diffraction}
在独立原子、运动学近似中，反射振幅为
\[
 F(h)=\sum_j o_j f_j(h)e^{2\pi i h^\T x_j}\,T_j(h),
 \qquad I(h)\propto |F(h)|^2,
\]
其中 $T_j$ 是选定的温度因子。
\begin{proposition}[平移消光]
若原子图样在有限平移子群下不变，则反射振幅含该平移群特征标之和；
非平凡特征标的和为零，所以相应反射消光。
\end{proposition}
\begin{proof}
将求和按平移轨道分组，每个轨道公共系数乘
$\sum_t e^{2\pi i h^\T t}$。
平移置换与非平凡特征标的求和恒等式给出零。
\end{proof}
螺旋轴和滑移面的消光也需带相位的完整操作，不能只查点群。
调制结构在 Fourier 展开下出现 $h+Q^\T m$ 的卫星位置；
强度仍取决于散射因子和具体调制，而非仅由子群标签决定。
官方交互显示的具体散射近似、峰形和版本必须另行核对\citep{OfficialHelp}；
磁中子散射需磁矩垂直于散射矢量的分量，不能由上述标量式直接替代。

三维显示只把坐标、半径、键阈值与矢量绘制成图；
阈值键线不是化学键或力常数的证明。
本节说明理论对象与预期不变量，没有审计 ISOVIZ 渲染器或衍射程序源代码。
对称模式也不是完整声子本征模：后者还需要质量和力常数；
SMODES 的对称分块不能自动给出频率\citep{Frozen1995,SmodesManual}。

\section{CIF、模式文件与精修接口}\label{sec:formats}
CIF 是数据语法和字典合同，不是算法正确性认证。
结构的晶格、对称操作、坐标和占位须相互一致；
msCIF/magCIF 另规定调制和磁量的语义\citep{CifSpec,CifCore,CifMs,CifMag}。
官方帮助把多数 \code{_iso_*} 模式标签描述为软件专属扩展
\citep{OfficialHelp}，不能凭字段名前缀宣称它们都是 IUCr 标准字段。

\begin{proposition}[可逆序列化的必要合同]
若写入器 $E$、读取器 $D$ 满足 $D(E(s))=s$（允许声明的坐标/排序等价），
则结构 $s$ 的语义在往返中保持。
仅满足语法正确、读取器不报错或文件哈希相等均不足以推出这个等式。
\end{proposition}
\begin{proof}
第一句就是左逆的定义。对反向推论，取遗漏占位或误写原点的语法合法文件，
即可得到解析成功而恢复结构不同的反例。
\end{proof}
TOPAS、FullProf、CIF 和 IsoVIZ 应保留相同的模式矩阵、幅度单位、列身份与应变顺序。
模式精修实际上把 $u=B a$ 代入实验模型，减少或重组参数；
能否改善拟合取决于实验、约束和可辨识性，不由换基单独保证
\citep{Refinement2007,Kerman2012}。
本地 \code{ISOVIZ_INPUT} 的 CSV 到 \code{amp} 写入必须沿用原始系数合同，
不能把显示用 normalized amplitude 当成任意模式文件的参数。

\section{能力覆盖、证据缺口与结论}
本次按官方帮助将能力分为：输入/识别、晶格与 setting、k/irrep、
物理表示与投影、Method 1--4、超空间、幅度与应变、
结构生成、畴、主序参量/不变量、子群树、衍射、可视化、导出/精修。
每类均有本文件相应解释以及 \code{coverage.json} 的来源、证明位置和本地模块映射。
这表示\emph{主题覆盖}，不是每项实现已被形式化证明。

仍不能无条件认证的项目包括：外部二进制与数据表的全部正确性；
非公度微观列投影与官网标签的完整对应；极端斜晶格中的版本特定有限应变规范；
容差改变空间群或原子匹配时的行为；所有 GUI 和衍射近似；所有输出格式的实际互操作性。
这些缺口没有被写成已完成的证明。

\paragraph{复用规则。}
修改算法时先定位对应的数学合同，声明输入域与约定，核查引用和反例；
再用不变量和独立数据验收实现。引用当前文献不能继承旧源码签名的测试结论。
任何新增功能必须同时补充文献/证明、实现适用范围与验收证据。
本文给出可供该流程使用的理论基础，而不替代独立同行复核、实验验证或源码审计。

\bibliographystyle{unsrtnat}
\bibliography{references}
\end{document}
